\documentclass{article}
\usepackage[utf8]{inputenc}
\usepackage[portuguese]{babel}
\usepackage{graphicx} 
\usepackage{fancyhdr} 
\usepackage[margin=2.5cm, top=4.5cm]{geometry}
\usepackage{tcolorbox}

% --- Configuração do Cabeçalho ---
\pagestyle{fancy}
\fancyhf{} 
\renewcommand{\headrulewidth}{0pt} 

\fancyhead[C]{%
    \begin{tcolorbox}[colback=white, colframe=black, boxrule=0.5pt, sharp corners, arc=0mm, boxsep=2pt, left=5pt, right=5pt, top=5pt, bottom=5pt]
        \begin{minipage}[c]{0.65\textwidth}
            \small 
            Empresa: << header.company|teX >> \\
            Data: << header.doc_date|teX >> \\
            Interessado: << header.interested|teX >> \\
            Referência: << header.reference|teX >> \\
            fck(MPa): << mix.fck28 or '-' >> \\
            Slump(cm): << mix.target_slump or '-' >>
        \end{minipage}%
        \begin{minipage}[c]{0.35\textwidth}
            \flushright 
            <% if header.logo_filename %>
            \includegraphics[height=2.5cm, keepaspectratio]{<< header.logo_filename >>}
            <% endif %>
        \end{minipage}
    \end{tcolorbox}
}

\setlength{\headheight}{95pt} 

\begin{document}

\thispagestyle{fancy} 
\begin{center}
    \textbf{\Large VALIDAÇÃO DE TRAÇO QUANTITATIVO}
\end{center}

\vspace{0.5cm}

\begin{table}[h]
    \centering
    \caption{Traço e Materiais}
    \vspace{0.2cm}
    
    \begin{tabular}{|c|c|c|c|c|}
        \hline 
        Material & Tipo & Traço (g/g) & Consumo (kg/m³) & Massa Específica (kg/m³) \\ 
        \hline
        <% for mat in mix.materials %>
        << mat.commercial_name|teX >> & << mat.material_type|teX >> & << mat.ratio >> & << mat.consumption >> & << mat.specific_mass_kg_m3 >> \\
        \hline
        <% endfor %>
    \end{tabular}
    
    \label{tab:meu_exemplo}
\end{table}

\vspace{0.3cm}
\noindent Relação água/cimento: << mix.water_cement_ratio or 'N/A' >>

\vspace{0.5cm}

\begin{table}[h] 
    \centering  
    \caption{Informações adicionais do traço} 
    \vspace{0.2cm}

    \begin{tabular}{|c|c|}
        \hline 
        Teor de Argamassa ($\alpha$) & << mix.alpha or '-' >> \\
        Massa especifica teórica (kg/m³) & << mix.theoretical_density or '-' >> \\
        Massa especifica real (kg/m³) & << mix.specific_mass or '-' >> \\
        Ar incorporado (m³/m³) & << mix.air_entrained or '-' >> \\
        \hline
    \end{tabular}
\end{table}

\end{document}